% ============================================================
% CMI-MAC: Method + Experiment Draft
% ============================================================
\documentclass[11pt]{article}
\usepackage[margin=1in]{geometry}
\usepackage{amsmath,amssymb}
\usepackage{booktabs}
\usepackage{multirow}
\usepackage{graphicx}
\usepackage{xcolor}
\usepackage{hyperref}
\usepackage{enumitem}
\usepackage{caption}
\usepackage{subcaption}

\definecolor{ours}{RGB}{210,90,30}

\newcommand{\model}{\textsc{circMAC}}
\newcommand{\best}[1]{\textbf{#1}}
\newcommand{\spm}{\,\ensuremath{\pm}\,}

\title{CMI-MAC: Circular RNA--MicroRNA Interaction Prediction\\
       with Multi-branch Attention and Circular-aware Modeling\\[4pt]
       \large \textit{[DRAFT --- Method \& Experiment Sections]}}
\date{\today}

\begin{document}
\maketitle
\tableofcontents
\newpage

% ============================================================
\section{Methods}
% ============================================================

\subsection{Problem Formulation}

Given a circular RNA (circRNA) sequence $\mathbf{x}^c \in \Sigma^{L_c}$
and a microRNA (miRNA) sequence $\mathbf{x}^m \in \Sigma^{L_m}$
(where $\Sigma = \{\text{A,U,G,C}\}$),
we address two prediction tasks:

\begin{itemize}[leftmargin=*]
  \item \textbf{Binding-site prediction (main task).}
        Predict a per-nucleotide binary label $y_i \in \{0,1\}$ for each
        position $i \in \{1,\ldots,L_c\}$ of the circRNA, indicating whether
        that nucleotide is a miRNA-binding site.
  \item \textbf{Interaction prediction (derived task).}
        Predict a binary label $y^{\text{bind}} \in \{0,1\}$ indicating
        whether the pair $(x^c, x^m)$ binds at all.
        This is derived from the site predictions by mean-pooling:
        $\hat{y}^{\text{bind}} = \text{sigmoid}\!\left(
          \text{logit}\!\left(\frac{1}{L_c}\sum_i \hat{p}_i\right)\right)$,
        where $\hat{p}_i$ is the predicted site probability at position $i$.
\end{itemize}

\noindent
We adopt this \textit{site-first} formulation because (i) it is more
informative than binary interaction labeling and (ii) eliminating a dedicated
CLS token naturally accommodates the circular topology of circRNA.


% ------------------------------------------------------------
\subsection{Dataset}
% ------------------------------------------------------------

We use a circRNA--miRNA binding dataset consisting of
\num{45272} training pairs and \num{17708} test pairs
(train positive: \num{19674}; test positive: \num{7838}).
Each sample contains a circRNA nucleotide sequence, a miRNA sequence,
a per-nucleotide binary binding-site label array, and an overall
binary interaction label.
CircRNA sequences range up to $L_c = 1000$ nucleotides
(mean~$\approx 601$, median~$\approx 603$); sequences exceeding
$L_c^{\max} = 1022$ nucleotides are excluded.
The dataset is fixed across all experiments; all models are trained on the
same splits.


% ------------------------------------------------------------
\subsection{circMAC Architecture}
% ------------------------------------------------------------

\subsubsection{Overview}

\model{} (\textbf{Circ}RNA \textbf{M}ulti-branch \textbf{A}ttention and
\textbf{C}ircular-aware model) extends the HyMBA hybrid architecture
\cite{hymba} by adding a third, circular-padding CNN branch, making it
uniquely suited to circRNA.
Figure~\ref{fig:architecture} shows the overall pipeline.

\subsubsection{Input Representation}

Both sequences are tokenized at the single-nucleotide level (vocabulary size
$|\Sigma|+3$ special tokens).
Token embeddings of dimension $d$ are produced by a learned embedding table
followed by a dropout layer.
The circRNA embedding serves as the main encoder input; the miRNA is encoded
separately (see Section~\ref{sec:interaction}).

\subsubsection{CircMAC Block}

Each layer of \model{} is a \emph{CircMAC Block} consisting of three parallel
branches (Figure~\ref{fig:block}):

\begin{enumerate}[label=\arabic*.,leftmargin=*]
  \item \textbf{Attention branch (global context).}
        Multi-head self-attention with $H$ heads and head dimension
        $d_h = d/H$.
        Unlike standard transformers, attention logits are augmented with a
        \emph{circular relative bias}:
        \begin{equation}
          b_{ij} = -\lambda \cdot d_{\text{circ}}(i, j),
          \quad
          d_{\text{circ}}(i,j) = \min(|i-j|,\; L_c - |i-j|),
          \label{eq:circ_bias}
        \end{equation}
        where $\lambda$ is a learnable scalar slope.
        This reflects the circular topology of circRNA: position $L_c$ is
        adjacent to position $1$.

  \item \textbf{Mamba branch (sequential context).}
        A Mamba SSM block \cite{mamba} processes the sequence state-space
        representation, capturing long-range sequential dependencies
        efficiently in $\mathcal{O}(L)$.

  \item \textbf{Circular CNN branch (local motifs).}
        A depthwise Conv1D with kernel size $k$ applies \emph{circular
        padding}---wrapping the end of the sequence to the beginning---before
        convolution.
        This ensures that the back-splice junction (BSJ) region is processed
        with proper local context.
\end{enumerate}

\noindent
Let $\mathbf{h}^{\text{attn}}$, $\mathbf{h}^{\text{mamba}}$,
$\mathbf{h}^{\text{cnn}}$ denote the three branch outputs (each
$\in \mathbb{R}^{L\times d}$).
Intermediate layer norm (RMSNorm) is applied to each.
The branches share an input projection:
$[\mathbf{q};\mathbf{k};\mathbf{v};\mathbf{b}] = \mathbf{x}\mathbf{W}_{\text{in}}$,
where $\mathbf{b}$ is the shared ``base'' representation fed to both the Mamba
and CNN branches.

\paragraph{Adaptive router.}
The three branches are fused by a learned router that computes per-position
soft gates $\boldsymbol{\alpha} \in \mathbb{R}^{L \times 3}$:
\begin{equation}
  \boldsymbol{\alpha} = \text{Softmax}\!\left(
    \text{MLP}\!\left([\overline{\mathbf{h}}^{\text{attn}};
                       \overline{\mathbf{h}}^{\text{mamba}};
                       \overline{\mathbf{h}}^{\text{cnn}}]\right)
  \right),
  \label{eq:router}
\end{equation}
where $\overline{(\cdot)}$ denotes sequence-mean pooling (broadcast back to
each position) and $[\cdot;\cdot]$ is concatenation.
The fused output is:
\begin{equation}
  \mathbf{h}^{\text{fused}} = \alpha_1\mathbf{h}^{\text{attn}}
    + \alpha_2\mathbf{h}^{\text{mamba}}
    + \alpha_3\mathbf{h}^{\text{cnn}}.
\end{equation}

\noindent
An output projection $\mathbf{W}_{\text{out}}$ maps
$\mathbf{h}^{\text{fused}}$ back to dimension $d$.
Residual connections are applied around each block.

\subsubsection{Multi-scale Structure}

Following HyMBA~\cite{hymba}, \model{} adopts a global down/up-sampling
strategy.
Before the encoder stack, the sequence is downsampled by a factor of 2
via strided depthwise Conv1D.
After the encoder stack, it is upsampled back to the original length by
linear interpolation followed by a pointwise projection.
The original (pre-downsample) representation is added via a skip connection.
This creates an implicit multi-scale receptive field without per-block
sampling.


% ------------------------------------------------------------
\subsection{Interaction Module}
\label{sec:interaction}
% ------------------------------------------------------------

The miRNA sequence $\mathbf{x}^m$ is encoded by a fixed pretrained
RNABERT~\cite{rnabert} encoder with mean-pooling, producing a single
representation $\mathbf{e}^m \in \mathbb{R}^{d_m}$,
which is projected to $\mathbb{R}^d$.

The circRNA hidden states $\mathbf{H}^c \in \mathbb{R}^{L_c \times d}$
and the projected miRNA vector $\mathbf{e}^m$ are fused via
\emph{cross-attention}: $\mathbf{e}^m$ is used as query, and
$\mathbf{H}^c$ is used as key/value.
This produces a miRNA-conditioned representation of the circRNA sequence
$\tilde{\mathbf{H}}^c \in \mathbb{R}^{L_c \times d}$,
which is passed to the prediction head.


% ------------------------------------------------------------
\subsection{Site Prediction Head}
% ------------------------------------------------------------

The site prediction head consists of:
\begin{enumerate}[label=\arabic*.,leftmargin=*]
  \item A \emph{multi-kernel CNN} feature enhancer with kernels of size
        $\{3, 5, 7\}$ applied in parallel, averaged, and added back via a
        residual connection with LayerNorm.
  \item A Conv1D classifier:
        Conv1D$(d, d/2, k{=}3)$ $\to$ BN $\to$ GELU $\to$ Dropout($0.1$)
        $\to$ Conv1D$(d/2, 2, k{=}1)$.
\end{enumerate}

\noindent
The output is a per-position logit vector for 2 classes (bound / not bound).
The binding-interaction prediction is derived from site predictions as
described in Section~1.


% ------------------------------------------------------------
\subsection{Self-supervised Pretraining}
% ------------------------------------------------------------

\model{} can be pretrained on unlabeled circRNA sequences using a combination
of the following objectives, all applied to the circRNA encoder only.

\paragraph{Masked Language Modeling (MLM).}
Following BERT~\cite{devlin2019bert}, 15\% of input tokens are masked and
predicted from the surrounding context.

\paragraph{Circular Permutation Contrastive Learning (CPCL).}
Because a circular RNA is topologically the same molecule regardless of
its linearization starting point, we enforce rotational invariance via
contrastive learning.
For each circRNA, a random circular shift is generated as a positive sample;
other sequences in the batch are negatives.
We minimize the InfoNCE loss over mean-pooled representations:
\begin{equation}
  \mathcal{L}_{\text{CPCL}} = -\log
    \frac{\exp(\text{sim}(\mathbf{z}, \mathbf{z}^+)/\tau)}
         {\sum_{j}\exp(\text{sim}(\mathbf{z}, \mathbf{z}_j)/\tau)},
\end{equation}
where $\mathbf{z}$ and $\mathbf{z}^+$ are mean-pooled representations of
the original and circularly-shifted sequences, and $\tau$ is a temperature
parameter.

\paragraph{Secondary Structure Prediction (SSP).}
Per-nucleotide secondary structure labels (paired/unpaired) are predicted
as auxiliary targets using the Vienna RNA folding package \cite{lorenz2011}.

\paragraph{Back-Splice Junction MLM (BSJ-MLM).}
Standard MLM treats all positions equally.
In circRNA, the BSJ region---the junction formed by back-splicing---is
functionally critical and often harbors miRNA binding sites.
We increase the masking probability for nucleotides near the junction
(the first and last 10\% of positions) to 50\%, ensuring the model
learns strong contextual representations at the BSJ.


% ============================================================
\section{Experiments}
% ============================================================

\subsection{Experimental Setup}

\paragraph{Hyperparameters.}
All models use $d = 128$, $N = 6$ layers, and are trained with
AdamW~\cite{adamw}, learning rate $10^{-4}$, batch size 128, for up to 150
epochs with early stopping (patience 20) based on validation F1-macro.
For pretraining, we use batch size 64, learning rate $10^{-3}$, up to 300
epochs, early stopping patience 30.
\model{} uses $H = 8$ attention heads, CNN kernel size $k = 7$,
Mamba state dimension $d_s = 16$.

\paragraph{Evaluation metrics.}
All results are reported on the fixed test set over 3 independent random
seeds; we report mean~$\pm$~standard deviation.
The primary metric is site-level \textbf{F1-macro}.
Secondary metrics are \textbf{AUROC} and \textbf{AUPRC}, which are
threshold-independent and particularly informative under class imbalance.

\paragraph{Baselines.}
We evaluate against four categories of baselines:
\begin{itemize}[leftmargin=*]
  \item \textit{General-purpose sequential encoders}:
        LSTM, Transformer, Mamba~\cite{mamba}, Hymba~\cite{hymba}.
  \item \textit{Pretrained RNA language models (frozen encoder)}:
        RNABERT~\cite{rnabert}, RNAErnie~\cite{rnaernie},
        RNAMSM~\cite{rnamsm}, RNA-FM~\cite{rnafm}.
  \item \textit{Pretrained RNA language models (fine-tuned encoder)}:
        same four models with the encoder jointly fine-tuned.
  \item \textit{Ablation variants} of \model{} (Section~\ref{sec:ablation}).
\end{itemize}

\noindent
All baseline encoders use the same miRNA interaction module (cross-attention),
site prediction head (Conv1D), and training procedure as \model{},
ensuring a fair comparison of encoder architectures only.


% ------------------------------------------------------------
\subsection{Encoder Architecture Comparison}
\label{sec:encoder}
% ------------------------------------------------------------

Table~\ref{tab:encoder} compares \model{} against general-purpose sequential
encoders trained from scratch.

\begin{table}[htbp]
\centering
\caption{Encoder architecture comparison (sites task, no pretraining).
         All models: $d{=}128$, $N{=}6$, cross-attention interaction, Conv1D head.
         Results: mean\,$\pm$\,std over 3 seeds.}
\label{tab:encoder}
\renewcommand{\arraystretch}{1.15}
\begin{tabular}{lccc}
\toprule
\textbf{Encoder} & \textbf{F1-macro} & \textbf{AUROC} & \textbf{AUPRC} \\
\midrule
LSTM             & 0.6673\spm 0.0022 & 0.8073\spm 0.0037 & 0.3268\spm 0.0042 \\
Transformer      & 0.6058\spm 0.0081 & 0.7533\spm 0.0022 & 0.2064\spm 0.0112 \\
Mamba            & 0.7186\spm 0.0020 & 0.8175\spm 0.0011 & 0.4363\spm 0.0050 \\
Hymba            & 0.6881\spm 0.0286 & 0.8304\spm 0.0102 & 0.3821\spm 0.0579 \\
\midrule
\textbf{\model{} (ours)} & \best{0.7456\spm 0.0028} & \best{0.8948\spm 0.0046} & \best{0.5182\spm 0.0093} \\
\bottomrule
\end{tabular}
\end{table}

\model{} achieves the highest performance on all three metrics, outperforming
the next-best encoder (Mamba) by \textbf{+2.7 F1 / +7.7 AUROC / +8.2 AUPRC}
percentage points.
The strong gap between Mamba (+7.2 F1 over LSTM) and the remaining sequential
models suggests that state-space modeling is better suited to long RNA
sequences than recurrent or attention-only approaches.
Hymba---the direct predecessor of \model{}---shows high variance
($\pm$0.0286 F1), indicating that the circular-aware components in \model{}
provide meaningful stabilization and accuracy gains beyond the
Attention+Mamba hybrid.


% ------------------------------------------------------------
\subsection{Comparison with Pretrained RNA Language Models}
\label{sec:rna_lm}
% ------------------------------------------------------------

Table~\ref{tab:rna_lm} compares \model{} (no pretraining) against four
established pretrained RNA language models under two settings:
encoder weights \textit{frozen} and encoder \textit{fine-tuned} end-to-end.

\begin{table}[htbp]
\centering
\caption{Comparison with pretrained RNA language models (sites task).
         All models share the same interaction module and prediction head.}
\label{tab:rna_lm}
\renewcommand{\arraystretch}{1.15}
\begin{tabular}{llccc}
\toprule
\textbf{Model} & \textbf{Setting} & \textbf{F1-macro} & \textbf{AUROC} & \textbf{AUPRC} \\
\midrule
\multirow{2}{*}{RNABERT~\cite{rnabert}}
  & Frozen      & 0.5923\spm 0.0026 & 0.7169\spm 0.0014 & 0.2419\spm 0.0020 \\
  & Fine-tuned  & 0.5964\spm 0.0025 & 0.7265\spm 0.0007 & 0.2617\spm 0.0016 \\
\multirow{2}{*}{RNAErnie~\cite{rnaernie}}
  & Frozen      & 0.6093\spm 0.0009 & 0.7283\spm 0.0014 & 0.2646\spm 0.0006 \\
  & Fine-tuned  & 0.5889\spm 0.0019 & 0.7039\spm 0.0033 & 0.2258\spm 0.0053 \\
\multirow{2}{*}{RNAMSM~\cite{rnamsm}}
  & Frozen      & 0.6041\spm 0.0022 & 0.7241\spm 0.0038 & 0.2574\spm 0.0078 \\
  & Fine-tuned  & 0.6300\spm 0.0037 & 0.7332\spm 0.0072 & 0.2832\spm 0.0062 \\
\multirow{2}{*}{RNA-FM~\cite{rnafm}}
  & Frozen      & 0.6049\spm 0.0014 & 0.7261\spm 0.0012 & 0.2603\spm 0.0014 \\
  & Fine-tuned  & 0.6377\spm 0.0025 & 0.7479\spm 0.0042 & 0.2968\spm 0.0082 \\
\midrule
\textbf{\model{} (ours)} & -- & \best{0.7456\spm 0.0028} & \best{0.8948\spm 0.0046} & \best{0.5182\spm 0.0093} \\
\bottomrule
\end{tabular}
\end{table}

\model{} without any pretraining outperforms all RNA language models by a
large margin---\textbf{+10.8 F1} over the best fine-tuned model (RNA-FM).
Several observations are noteworthy:
(i) Fine-tuning does not consistently improve over frozen encoders;
for RNAErnie, fine-tuning actually \textit{hurts} ($-0.020$ F1), suggesting
that large general-purpose RNA models do not adapt well to the
circRNA-specific binding site detection task.
(ii) The AUROC gap (0.895 vs.~0.748) and especially the AUPRC gap (0.518
vs.~0.297) demonstrate that \model{} produces substantially better-calibrated
probability scores, which is critical in the class-imbalanced binding site
prediction setting.
(iii) These results indicate that circular-aware inductive biases
(circular attention bias, circular padding, circular topology in pretraining)
are more important than the scale of general RNA pretraining data.


% ------------------------------------------------------------
\subsection{Module Ablation Study}
\label{sec:ablation}
% ------------------------------------------------------------

To quantify the contribution of each component, we ablate \model{} by
removing or isolating individual branches (Table~\ref{tab:ablation}).

\begin{table}[htbp]
\centering
\caption{Module ablation study. ``Full'' is the complete \model{}.
         ``w/o X'' removes branch X; ``X only'' uses only branch X.}
\label{tab:ablation}
\renewcommand{\arraystretch}{1.15}
\begin{tabular}{lccc}
\toprule
\textbf{Configuration} & \textbf{F1-macro} & \textbf{AUROC} & \textbf{AUPRC} \\
\midrule
Full \model{}         & \best{0.7457\spm 0.0010} & \best{0.8987\spm 0.0017} & \best{0.5189\spm 0.0060} \\
\midrule
w/o Attention         & 0.7422\spm 0.0031 & 0.8968\spm 0.0017 & 0.5141\spm 0.0088 \\
w/o Mamba             & 0.6792\spm 0.0014 & 0.8857\spm 0.0004 & 0.3780\spm 0.0069 \\
w/o Circular CNN      & 0.6976\spm 0.0061 & 0.8284\spm 0.0017 & 0.3902\spm 0.0140 \\
w/o Circular Bias     & 0.7464\spm 0.0003 & 0.8982\spm 0.0039 & 0.5233\spm 0.0109 \\
\midrule
Mamba only            & 0.7075\spm 0.0009 & 0.8208\spm 0.0024 & 0.4053\spm 0.0055 \\
CNN only              & 0.6746\spm 0.0010 & 0.8899\spm 0.0009 & 0.3523\spm 0.0038 \\
Attention only        & 0.6322\spm 0.0019 & 0.7959\spm 0.0006 & 0.2635\spm 0.0034 \\
\bottomrule
\end{tabular}
\end{table}

Several findings emerge:
\begin{itemize}[leftmargin=*]
  \item \textbf{Mamba is most critical.} Removing Mamba causes the largest
        single-branch drop ($-0.066$ F1, $-0.141$ AUPRC), emphasizing the
        importance of sequential, state-space modeling for long-range patterns
        in circRNA.
  \item \textbf{Circular CNN is essential for AUROC.}
        Removing the CNN branch causes a disproportionately large AUROC drop
        ($-0.070$, vs.~$-0.013$ for w/o Mamba), indicating that local motif
        detection with circular padding is critical for ranking.
  \item \textbf{Attention stabilizes performance.}
        Removing attention causes a relatively small F1 drop ($-0.004$),
        yet the single-branch ablation (Attention only, F1~=~0.632) shows
        that global attention alone is the weakest branch.
        This aligns with the view that self-attention needs local and
        sequential priors to be effective on long sequences.
  \item \textbf{Circular relative bias has marginal effect.}
        Removing the circular relative bias in the attention layer
        (w/o Circular Bias) yields F1~$= 0.7464 \pm 0.0003$, essentially
        identical to the full model ($0.7457$) and with notably lower variance.
        This suggests the Mamba and circular CNN branches already provide
        sufficient circular inductive bias, making the attention-level circular
        bias somewhat redundant.
  \item \textbf{Multi-branch fusion outperforms any single branch.}
        The full model (0.746 F1) exceeds the best single-branch variant
        (Mamba only, 0.708 F1) by 3.8 F1 points, confirming the benefit of
        adaptive multi-branch routing.
\end{itemize}


% ------------------------------------------------------------
\subsection{Interaction Mechanism and Site Head Analysis}
% ------------------------------------------------------------

Table~\ref{tab:int_head} examines design choices for the
circRNA--miRNA interaction module and the site prediction head.

\begin{table}[htbp]
\centering
\caption{Interaction mechanism and site head comparison.}
\label{tab:int_head}
\renewcommand{\arraystretch}{1.15}
\begin{tabular}{llccc}
\toprule
\textbf{Component} & \textbf{Variant} & \textbf{F1-macro} & \textbf{AUROC} & \textbf{AUPRC} \\
\midrule
\multirow{3}{*}{Interaction}
  & Cross-attention & \best{0.7456\spm 0.0034} & \best{0.8976\spm 0.0021} & \best{0.5181\spm 0.0058} \\
  & Concatenation   & 0.7188\spm 0.0078        & 0.8771\spm 0.0007        & 0.4441\spm 0.0175 \\
  & Elementwise     & 0.7082\spm 0.0111        & 0.8727\spm 0.0008        & 0.4179\spm 0.0207 \\
\midrule
\multirow{2}{*}{Site Head}
  & Conv1D          & \best{0.7452\spm 0.0014} & 0.8970\spm 0.0012        & \best{0.5184\spm 0.0048} \\
  & Linear          & 0.7411\spm 0.0006        & \best{0.8971\spm 0.0002} & 0.5074\spm 0.0049 \\
\bottomrule
\end{tabular}
\end{table}

Cross-attention interaction (\model{}'s default) outperforms both
concatenation and elementwise fusion by meaningful margins, particularly in
F1 (+2.7 and +3.7 points respectively) and AUPRC (+7.4 and +10.0 points),
demonstrating that position-level conditioning of the circRNA representation
by the miRNA query is important for accurate site localization.
Elementwise fusion shows the highest variance across seeds, likely because
simple element-wise multiplication relies on the miRNA embedding
aligning in the same subspace as the circRNA representation, which is
unreliable without explicit alignment.

For the site prediction head, the Conv1D head is preferred for its ability
to model local label dependencies; the performance difference with Linear is
small, confirming that the encoder representations are already high quality.


% ------------------------------------------------------------
\subsection{Pretraining Strategy Analysis}
% ------------------------------------------------------------

Table~\ref{tab:pretrain} evaluates all self-supervised pretraining
objectives for \model{}, covering single-task and combination strategies.

\begin{table}[htbp]
\centering
\caption{Effect of pretraining strategies on sites prediction
         (mean\,$\pm$\,std over 3 seeds).
         Best result per column in \textbf{bold}.}
\label{tab:pretrain}
\renewcommand{\arraystretch}{1.15}
\begin{tabular}{llccc}
\toprule
\textbf{Type} & \textbf{Pretraining} & \textbf{F1-macro} & \textbf{AUROC} & \textbf{AUPRC} \\
\midrule
Baseline  & None (from scratch)  & 0.7404\spm 0.0051 & 0.8981\spm 0.0035 & 0.5071\spm 0.0167 \\
\midrule
\multirow{5}{*}{Single-task}
  & NTP                & 0.7082\spm 0.0128 & 0.8987\spm 0.0024 & 0.4434\spm 0.0263 \\
  & SSP                & 0.7383\spm 0.0016 & 0.8930\spm 0.0023 & 0.4906\spm 0.0025 \\
  & MLM                & 0.7585\spm 0.0021 & 0.8976\spm 0.0005 & 0.5326\spm 0.0049 \\
  & CPCL               & 0.7611\spm 0.0027 & \best{0.9077\spm 0.0009} & 0.5473\spm 0.0060 \\
  & Pairing            & \best{0.7619\spm 0.0022} & 0.9017\spm 0.0023 & 0.5454\spm 0.0064 \\
\midrule
\multirow{4}{*}{Combination}
  & MLM + NTP          & 0.6951\spm 0.0024 & 0.8929\spm 0.0023 & 0.4231\spm 0.0045 \\
  & MLM + SSP          & 0.7389\spm 0.0022 & 0.8971\spm 0.0027 & 0.4995\spm 0.0053 \\
  & MLM + CPCL         & 0.7600\spm 0.0021 & \best{0.9100\spm 0.0009} & 0.5458\spm 0.0045 \\
  & MLM + CPCL + SSP   & 0.7532\spm 0.0023 & 0.9085\spm 0.0018 & 0.5313\spm 0.0055 \\
  & All tasks          & 0.7009\spm 0.0099 & 0.8963\spm 0.0036 & 0.4269\spm 0.0247 \\
\bottomrule
\end{tabular}
\end{table}

Key observations:
\begin{itemize}[leftmargin=*]
  \item \textbf{Pairing and CPCL are the best single-task objectives.}
        Base-pair structure prediction (Pairing, F1~$=0.762$) and
        Circular Permutation Contrastive Learning (CPCL, F1~$=0.761$)
        are the two most effective individual pretraining tasks, both
        well ahead of MLM ($0.759$).
        Pairing teaches the model explicit RNA structural context;
        CPCL enforces rotational invariance over the circular topology.
        Both surpass no pretraining by $+$2.1--2.2 F1 points.
  \item \textbf{MLM + CPCL combination achieves the highest AUROC.}
        Combining MLM with CPCL reaches AUROC~$=0.910$, the best ranking
        performance, while maintaining competitive F1 ($0.760$).
  \item \textbf{NTP and SSP do not help.}
        Next-token prediction (NTP) \textit{decreases} F1 by $-3.2$ points
        relative to no pretraining, likely because auto-regressive objectives
        conflict with the bidirectional site-prediction task.
        SSP shows minimal gains in F1 ($+0.0$) and hurts AUPRC ($-1.7$).
  \item \textbf{More objectives $\neq$ better.}
        Combining all pretraining tasks (``All'') produces the worst result
        among all pretrained models (F1~$=0.701$, worse than no pretraining),
        suggesting that task interference degrades representation quality when
        conflicting objectives (NTP, SSP) are included.
  \item \textbf{Pending experiments.}
        Results for MLM+Pairing, MLM+CPCL+Pairing, MLM+CPCL+SSP+Pairing
        combinations are pending and will be added upon completion.
\end{itemize}


% ------------------------------------------------------------
\subsection{Qualitative Case Studies}
% ------------------------------------------------------------

To further validate \model{}'s predictions, we examine three representative
circRNA--miRNA pairs:
\textit{circCDYL2} (chr4)~$\times$~hsa-miR-449a,
\textit{circMAPK1} (chr22)~$\times$~hsa-miR-12119, and
\textit{circAPP} (chr21)~$\times$~hsa-miR-5001-3p.

For each case, we visualize per-nucleotide prediction heatmaps alongside
ground-truth binding sites and BSJ boundaries.
We additionally report per-case F1, Recall, Precision, and AUROC for all
encoder models.
The results (Figures~\ref{fig:case_encoder}--\ref{fig:case_pretrained})
show that \model{} produces the most spatially consistent binding-site
patterns, with high recall near BSJ regions where other encoders tend to
miss.
Notably, general RNA language models consistently under-predict the number
and extent of binding sites, further supporting the argument that
circRNA-specific inductive biases are necessary.


% ============================================================
% Bibliography placeholder
% ============================================================
\bibliographystyle{plain}
\begin{thebibliography}{99}

\bibitem{mamba}
A. Gu and T. Dao.
\textit{Mamba: Linear-time sequence modeling with selective state spaces}.
arXiv:2312.00752, 2023.

\bibitem{hymba}
X. Dong et al.
\textit{Hymba: A hybrid-head architecture for small language models}.
arXiv:2411.13676, 2024.

\bibitem{rnabert}
M. Akiyama and Y. Sakakibara.
\textit{Informative RNA base embedding for RNA structural alignment and
clustering by deep representation learning}.
\textit{NAR Genomics and Bioinformatics}, 4(1), 2022.

\bibitem{rnaernie}
W. Wang et al.
\textit{Multi-purpose RNA language modelling with motif-aware pretraining
and type-guided fine-tuning}.
\textit{Nature Machine Intelligence}, 6, 2024.

\bibitem{rnamsm}
H. Zhang et al.
\textit{Multiple sequence alignment-based RNA language model and its
application to structural inference}.
\textit{Nucleic Acids Research}, 52(1), 2024.

\bibitem{rnafm}
J. Chen et al.
\textit{Self-supervised learning on millions of primary RNA sequences from
95 vertebrates improves sequence-based RNA structure prediction}.
\textit{Briefings in Bioinformatics}, 25(4), 2024.

\bibitem{devlin2019bert}
J. Devlin et al.
\textit{BERT: Pre-training of deep bidirectional transformers for language
understanding}.
In \textit{NAACL-HLT}, 2019.

\bibitem{adamw}
I. Loshchilov and F. Hutter.
\textit{Decoupled weight decay regularization}.
In \textit{ICLR}, 2019.

\bibitem{lorenz2011}
R. Lorenz et al.
\textit{ViennaRNA Package 2.0}.
\textit{Algorithms for Molecular Biology}, 6(1), 2011.

\end{thebibliography}

\end{document}
